% chapters/13/section-06.tex
\section{Kamera, Rasa Gerak, dan Penutup}
\label{sec:kamera-rasa}

\subsection{Kamera yang Mengikuti Pemain}

Kamera pada permainan platformer biasanya mengikuti pemain, dan cara
mengikutinya menentukan sebagian besar rasa gerak permainannya. Skripnya
singkat, tetapi satu hal perlu diperhatikan dengan saksama.

\begin{kodekeluaran}
void LateUpdate()
{
    Vector3 target = new Vector3(pemain.position.x, pemain.position.y,
                                 transform.position.z);
    transform.position = Vector3.Lerp(transform.position, target,
                                      followSpeed * Time.deltaTime);
}
\end{kodekeluaran}

Kamera dipindahkan pada fungsi \code{LateUpdate}, yaitu fungsi
yang dipanggil mesin \textbf{setelah} seluruh \code{Update} dan
\code{FixedUpdate} selesai. Urutan itu penting: kamera harus dipindahkan
berdasarkan letak pemain yang sudah diperbarui pada bingkai itu, bukan letak
pemain pada awal bingkai.

Perhatikan pula pemakaian \code{Vector3.Lerp}, yaitu pengambilan nilai antara
dua letak menurut satu angka. Cara itu membuat kamera menyusul pemain secara
berangsur, bukan mengikutinya dengan tepat. Nilai pengali yang kecil
menghasilkan kamera yang lebih lambat menyusul, dan itulah salah satu pengatur
rasa gerak yang paling berpengaruh.

\subsection{Mengatur Rasa Gerak}

Rasa gerak permainan platformer sebagian besar ditentukan oleh lima nilai, dan
seluruhnya dapat diatur tanpa mengubah kode.

\begin{itemize}
  \item \textbf{Kecepatan mendatar} menentukan seberapa gesit pemain
        berpindah.
  \item \textbf{Kecepatan lompatan} menentukan tinggi lompatan, sebagaimana
        dijelaskan pada Bagian~6.3.
  \item \textbf{\istilah{Gravity Scale}} menentukan seberapa cepat pemain
        jatuh kembali, dan bersama kecepatan lompatan menentukan lama pemain
        berada di udara.
  \item \textbf{Kelambatan kamera} menentukan seberapa jauh kamera tertinggal
        dari pemain.
  \item \textbf{Daerah pemeriksaan tanah} menentukan seberapa mudah pemain
        dapat melompat, khususnya di tepi bongkahan.
\end{itemize}

\begin{tipsbox}
Kelima nilai itu seluruhnya merupakan pilihan perancang, dan
tidak satu pun berasal dari perhitungan. Inilah yang dimaksud dengan
perancangan rasa gerak pada Bagian~6.3: yang menentukan rasa permainan bukan
tinggi lompatan itu sendiri, melainkan perbandingan antara nilai-nilai
tersebut.
\end{tipsbox}

\subsection{Mengakhiri Blok III}

Bab ini menutup \CPMK{3}, yaitu kemampuan menggunakan \istilah{game engine}
untuk membangun level. Tabel~\ref{tab:perkakas-unity} memuat perkakas yang
telah lengkap pada dua bab terakhir.

\begin{table}[htbp]
\centering
\caption{Perkakas Unity yang telah lengkap beserta bab yang membahasnya}
\label{tab:perkakas-unity}
{\setlength{\tabcolsep}{4pt}\footnotesize%
\begin{tabularx}{\textwidth}{@{}L{3.4cm} L{2.4cm} Y@{}}
\toprule
\headrow \textbf{Perkakas} & \textbf{Dibahas pada} & \textbf{Kegunaan} \\
\midrule
Antarmuka dan penyusunan benda & Bab~XII &
Mengelola project, scene, dan benda beserta komponennya. \\
\addlinespace
Skrip dan urutan pemanggilannya & Bab~XII &
Menempatkan pekerjaan pada fungsi yang tepat. \\
\addlinespace
Penyusunan level & Bab~XIII &
Menyusun ruang permainan secara \istilah{greyboxing} dan dengan
\istilah{tilemap}. \\
\addlinespace
Fisika dan tabrakan & Bab~XIII &
Mengonfigurasi \istilah{Collider}, \istilah{Rigidbody}, dan
\istilah{trigger}, serta memperbaiki masalahnya. \\
\bottomrule
\end{tabularx}}
\end{table}

Bab~XIV sampai Bab~XVI melengkapi level itu dengan aset visual, antarmuka,
audio, dan pengujian prototipe. Seluruhnya berpijak pada level yang baru saja
Anda bangun.
